\section{Arranque y configuracion del nodo}

\subsection{Build: fat jar con Maven}

El proyecto se compila con Maven y produce un \emph{fat jar} ejecutable
descrito en \codigo{pom.xml}. El artefacto es \codigo{tesoreria-distribuida}
version \codigo{1.0.0}, empaquetado como \codigo{jar} y compilado contra Java~17
(\codigo{maven.compiler.release} = 17). El comando habitual es
\codigo{mvn clean package}, que dispara el plugin
\codigo{maven-assembly-plugin} (version 3.7.1) con el descriptor
\codigo{jar-with-dependencies}: ese descriptor reune el codigo del nodo y todas
sus dependencias en un solo jar autocontenido cuyo manifiesto fija la clase
principal \codigo{mx.ipn.escom.tesoreria.app.Node}. Asi un unico jar arranca
indistintamente como lider o como replica, segun las variables de entorno.

\begin{lstlisting}[language=none,caption={Plugin assembly: fat jar y clase principal (pom.xml)}]
<artifactId>maven-assembly-plugin</artifactId>
<version>3.7.1</version>
<configuration>
    <descriptorRefs>
        <descriptorRef>jar-with-dependencies</descriptorRef>
    </descriptorRefs>
    <archive>
        <manifest>
            <mainClass>mx.ipn.escom.tesoreria.app.Node</mainClass>
        </manifest>
    </archive>
</configuration>
\end{lstlisting}

El sistema es Java puro, sin frameworks: la lista de dependencias declaradas se
limita a tres bibliotecas, \codigo{org.mindrot:jbcrypt:0.4} (hash de
contrasenas), \codigo{com.auth0:java-jwt:4.3.0} (emitir y validar JWT) y
\codigo{com.google.code.gson:gson:2.13.1} (parseo de JSON). Es importante para
la rubrica que el \codigo{pom.xml} \emph{no} traiga ningun driver JDBC ni base
de datos embebida: el requisito de usar un SGBD se cumple con cero, porque el
estado vive en memoria (el ledger) y la durabilidad se resuelve contra Cloud
Storage. De hecho el comentario del propio \codigo{pom.xml} aclara que no se
permite el SDK de Google Cloud, y que GCS se alcanza por su API REST JSON cruda
usando \codigo{java.net.http.HttpClient}.

\subsection{Punto de entrada: Node}

La clase \codigo{app/Node.java} es el punto de entrada del proceso. Su
\codigo{main()} recibe opcionalmente el puerto HTTP como \codigo{args[0]} (por
defecto 8080) y orquesta el arranque en una secuencia fija. Primero lee la
configuracion con \codigo{NodeConfig.fromEnv()}; despues carga el dataset de
cuentas en un \codigo{Ledger} nuevo con
\codigo{Dataset.loadInto(ledger, ...)} y, si la carga devuelve cero cuentas,
falla de inmediato con \codigo{IllegalStateException} en vez de servir un ledger
vacio (sintoma tipico de un CSV ausente o truncado). El ledger ya poblado se
envuelve en un \codigo{Bank}, que es quien posee la secuencia de \emph{commits}
y el reparto a los \codigo{CommitListener}.

\begin{lstlisting}[language=Java,caption={Inicio de Node.main(): config, dataset y bank}]
int port = (args.length > 0) ? Integer.parseInt(args[0]) : 8080;

NodeConfig config = NodeConfig.fromEnv();

Ledger ledger = new Ledger();
int loaded = Dataset.loadInto(ledger, Path.of(config.datasetPath()), config.idBase());
if (loaded <= 0) {
    throw new IllegalStateException("dataset loaded 0 accounts from "
            + config.datasetPath() + "; refusing to start with an empty ledger");
}
// ...
Bank bank = new Bank(ledger);
\end{lstlisting}

A continuacion el nodo se ramifica segun su rol, determinado por
\codigo{config.isLeader()}. Si es \textbf{lider} y hay Cloud Storage configurado
(\codigo{TES\_GCS\_KEYFILE} y \codigo{TES\_BUCKET}), construye un
\codigo{GcsStore} y un \codigo{Journal}, ejecuta la recuperacion en frio con
\codigo{journal.recover(bank::applyReplicated)} para reaplicar las
transferencias durables, registra el journal como \codigo{CommitListener} y
agrega un \emph{shutdown hook} que drena la cola durable al apagar. Tambien
levanta el \codigo{ReplicaFeed} en el puerto de replicacion y lo registra como
otro listener, de modo que cada commit en vivo se transmita por TCP a las
replicas conectadas. Si es \textbf{replica}, opcionalmente restaura su
\codigo{ReplicaCheckpoint} desde GCS \emph{antes} de iniciar la replicacion (para
que el \codigo{CATCHUP} arranque desde la marca exacta ya aplicada y no desde 0),
y luego un \codigo{ReplicaSync} se pone al dia desde el lider y aplica los
commits en vivo a traves del banco.

Por ultimo, \codigo{main()} arma la pila de seguridad
(\codigo{Tokens}, \codigo{CredentialStore}, \codigo{Authenticator}), construye el
proveedor de metricas \codigo{NodeStats}, registra la tabla de rutas con las
rutas exactas del contrato y arranca el servidor. La ultima linea bloquea el
hilo principal mientras los reactores y los workers atienden las peticiones.

\begin{lstlisting}[language=Java,caption={Tabla de rutas y arranque del servidor (Node.main())}]
Routes routes = new Routes();
routes.register("POST", "/api/register", new RegisterEndpoint(authenticator));
routes.register("POST", "/api/login", new LoginEndpoint(authenticator));
routes.register("GET", "/api/accounts/{id}", new BalanceEndpoint(authenticator, ledger));
routes.register("POST", "/api/transactions/transfer", new TransferEndpoint(authenticator, bank));
routes.register("GET", "/api/stats", new StatsEndpoint(stats));
routes.register("GET", "/panel", new PanelEndpoint(stats, config));
routes.register("GET", "/", new DashboardEndpoint());

Server server = new Server(port, routes, config.workers(), config.reactors());
server.start();

// Block the main thread while the IoLoop and workers serve requests.
Thread.currentThread().join();
\end{lstlisting}

El proveedor de metricas se instancia como
\codigo{NodeStats(ledger, bank, journal, config)} (ver \codigo{app/NodeStats.java}):
toma el numero de cuentas y el saldo total del \codigo{Ledger}, el conteo de
transferencias y el ultimo id del \codigo{Bank}, y el conteo en GCS del
\codigo{Journal} (que vale 0 en las replicas, donde el journal es null). Las
metricas de host (CPU, RAM, disco) se leen con Java crudo desde
\codigo{/proc/stat}, \codigo{/proc/meminfo} y el sistema de archivos, sin
bibliotecas externas.

\subsection{Configuracion por entorno (TES\_*)}

Toda la configuracion del nodo se inyecta por variables de entorno con prefijo
\codigo{TES\_}, leidas una sola vez en \codigo{NodeConfig.fromEnv()} dentro de
\codigo{app/NodeConfig.java}. El rol se deriva de \codigo{TES\_LEADER\_HOST}:
vacio o ausente significa que el proceso es el \textbf{lider} (nodo-1); cualquier
otro valor nombra al host del lider que esta \textbf{replica} debe seguir. El
secreto \codigo{TES\_JWT\_SECRET} debe ser el MISMO en los tres nodos para que un
token emitido en cualquiera valide en todos. La siguiente tabla resume cada
variable con su significado y su valor por defecto (los defaults vienen de las
constantes \codigo{DEFAULT\_*} de la clase).

\begin{center}
\begin{tabular}{|l|l|p{6.4cm}|}
\hline
\textbf{Variable} & \textbf{Default} & \textbf{Significado} \\
\hline
\codigo{TES\_DATASET} & (ninguno) & Ruta al CSV de cuentas que carga el nodo. \\
\codigo{TES\_JWT\_SECRET} & (ninguno) & Secreto HMAC para firmar/validar JWT; identico en los 3 nodos. \\
\codigo{TES\_NODE\_ID} & (ninguno) & Identificador legible del nodo, p.ej. \codigo{nodo-1}. \\
\codigo{TES\_PEERS} & \codigo{""} & URLs base de los pares, separadas por coma, para la agregacion del panel. \\
\codigo{TES\_LEADER\_HOST} & \codigo{""} & Host del lider a seguir; vacio = este nodo ES el lider. \\
\codigo{TES\_REPL\_PORT} & 9090 & Puerto TCP del feed de replicacion del lider. \\
\codigo{TES\_BUCKET} & (ninguno) & Bucket de Cloud Storage para el journal durable. \\
\codigo{TES\_GCS\_KEYFILE} & (ninguno) & Ruta al key file de la cuenta de servicio para autenticar contra GCS. \\
\codigo{TES\_WORKERS} & 16 & Tamano del pool de hilos que procesan peticiones. \\
\codigo{TES\_ID\_BASE} & 1 & Id asignado a la primera cuenta del dataset. \\
\codigo{TES\_REPLICA\_CHECKPOINT\_INTERVAL\_SECS} & 10 & Segundos entre flushes del checkpoint de la replica a GCS. \\
\codigo{TES\_REPLICA\_CHECKPOINT\_PATH} & (deriva) & Nombre del objeto del checkpoint; si es null usa \codigo{checkpoint/<nodeId>.json}. \\
\codigo{TES\_REACTORS} & 1 por CPU & Numero de reactores (selectores) del servidor HTTP. \\
\hline
\end{tabular}
\end{center}

La lectura usa dos ayudantes privados: \codigo{env(nombre, def)} devuelve el
default cuando la variable esta ausente o vacia, e \codigo{intEnv(nombre, def)}
parsea un entero y cae al default si la variable falta o no es numerica. El
metodo \codigo{isLeader()} centraliza la decision de rol: devuelve \codigo{true}
cuando \codigo{leaderHost} es null, vacio o solo espacios en blanco.

\begin{lstlisting}[language=Java,caption={Lectura de las variables TES\_* en NodeConfig.fromEnv()}]
public static NodeConfig fromEnv() {
    return new NodeConfig(
            env("TES_DATASET", null),
            env("TES_JWT_SECRET", null),
            env("TES_NODE_ID", null),
            env("TES_PEERS", ""),
            env("TES_LEADER_HOST", ""),
            intEnv("TES_REPL_PORT", DEFAULT_REPL_PORT),
            // ...
            intEnv("TES_WORKERS", DEFAULT_WORKERS),
            intEnv("TES_ID_BASE", DEFAULT_ID_BASE),
            // ...
            intEnv("TES_REACTORS", DEFAULT_REACTORS));
}
\end{lstlisting}

\subsection{Carga del dataset de cuentas}

Las cuentas se cargan desde el CSV con \codigo{Dataset.loadInto()} en
\codigo{core/Dataset.java}. El archivo de referencia es
\codigo{material-profesor/alumnos.csv}, entregado por el profesor y consumido tal
cual (no se regenera ni se edita). Es CSV plano de cuatro columnas y \emph{sin}
fila de encabezado, con el formato
\codigo{nombre,apellido1,apellido2,saldo}; por eso la primera linea tambien se
parsea, no se salta. Como no hay columna de id, el id de cada cuenta es su
posicion en el archivo contando desde \codigo{idBase} (normalmente 1, via
\codigo{TES\_ID\_BASE}): la primera linea de datos pasa a ser la cuenta
\codigo{idBase}, la segunda \codigo{idBase+1}, y asi sucesivamente. Que todos los
nodos lean el mismo archivo es justo lo que hace significativa la verificacion de
conservacion del dinero en el cluster.

El nombre del titular es la union con espacios de los tres campos de nombre. El
saldo se convierte de su texto decimal a centavos enteros con
\codigo{Money.toCents(String)}, de modo que ningun punto flotante toque jamas un
saldo: los importes se manejan internamente en centavos. Las lineas vacias o con
menos de cuatro columnas se descartan, y el metodo devuelve el numero de cuentas
efectivamente cargadas (el valor que \codigo{Node.main()} exige que sea mayor que
cero).

\begin{lstlisting}[language=Java,caption={Asignacion de id por indice de fila y saldo en centavos (Dataset.loadInto)}]
int id = idBase;
int loaded = 0;
// ...
String[] cols = line.split(",", -1);
if (cols.length < 4) {
    continue;
}
String owner = cols[0].trim() + " " + cols[1].trim() + " " + cols[2].trim();
long cents = Money.toCents(cols[3].trim());
ledger.add(new Account(id, owner, cents));
id++;
loaded++;
\end{lstlisting}
